% ============================================================================
% Unified Supplementary Information
% Emotional Dysregulation as the Core of Transdiagnostic Psychopathology
% Nature Reviews Psychology
% ============================================================================
% Standalone LaTeX content — no preamble, no \begin{document}.
% Requires: longtable, booktabs, hyperref, enumitem, geometry, xcolor,
%           array, makecell, amsmath packages in the master document.
% ============================================================================

\sloppy

\sectitle{Supplementary Information}

\bigskip

This Supplementary Information documents the computational search strategy,
complete tool call log (54 calls across 5 waves), statistical methods,
evidence grading for every major claim, limitations and potential biases,
author contributions, agent prompts, negative and null results, and data
availability for the accompanying review.  Every factual claim in the main
text is traceable to a specific database query recorded here.  No claims
rely on model knowledge alone; all are grounded in verifiable database
returns.  This appendix is designed so that an independent researcher can
retrace every analytical step from tool discovery through final inference.

% ============================================================================
% SI-1: SEARCH STRATEGY
% ============================================================================
\sectitle{SI-1: Search Strategy}

\subsection*{SI-1.1\quad Infrastructure}

All database queries were executed through ToolUniverse v1.1.5, an open-source
scientific MCP (Model Context Protocol) server that exposes approximately
1,700 tool wrappers spanning more than 600 scientific databases and APIs.
ToolUniverse operates in \emph{compact mode}, presenting five proxy tools to
the orchestrating language model:

\begin{enumerate}[nosep]
  \item \texttt{grep\_tools} --- keyword-based tool discovery (exact match).
  \item \texttt{find\_tools} --- natural-language / embedding-based tool
        discovery.
  \item \texttt{list\_tools} --- category browsing of the full tool catalogue.
  \item \texttt{get\_tool\_info} --- retrieval of the exact JSON parameter
        schema for a given tool (mandatory before execution).
  \item \texttt{execute\_tool} --- execution of any tool with a validated
        argument dictionary.
\end{enumerate}

\subsection*{SI-1.2\quad Multi-Strategy Tool Discovery}

Tool discovery used a three-tier pipeline:

\begin{itemize}[nosep]
  \item \textbf{Tool\_Finder\_LLM} (primary): an autonomous planning agent
        backed by Claude CLI (primary) with Ollama as a local fallback.
        Tool\_Finder\_LLM returned an empty result set on both invocations
        due to a known 120-second timeout under deep reasoning workloads
        over the 1,700-tool catalogue.  This failure is documented
        transparently (see SI-5) and triggered a fallback to keyword
        discovery.
  \item \textbf{grep\_tools}: keyword search against tool names and
        descriptions.  Used to locate domain-specific tools (e.g.,
        \texttt{gwas\_search\_studies}, \texttt{AllenBrain\_search\_genes}).
  \item \textbf{find\_tools}: embedding-based semantic search.  Used when
        keyword search returned no results or when the query was conceptual
        rather than lexical (e.g., ``neuroanatomy gene expression atlas'').
\end{itemize}

\subsection*{SI-1.3\quad Search Domains}

Ten search domains were defined \emph{a priori} based on the review's scope.
Each domain maps to one or more ToolUniverse-wrapped databases:

{\small
\begin{longtable}{@{}p{3.2cm}p{10.2cm}@{}}
\toprule
\textbf{Domain} & \textbf{Databases Queried} \\
\midrule
\endfirsthead
\toprule
\textbf{Domain} & \textbf{Databases Queried} \\
\midrule
\endhead
\bottomrule
\endlastfoot
Literature &
  PubMed, OpenAlex, Crossref, Europe PMC, Semantic Scholar \\
\midrule
Genomics &
  GWAS Catalog (NHGRI-EBI) \\
\midrule
Molecular biology &
  UniProt, STRING v12, Reactome, KEGG \\
\midrule
Pharmacology &
  PubChem, FAERS, DGIdb, PharmGKB, DrugBank \\
\midrule
Clinical trials &
  ClinicalTrials.gov \\
\midrule
Neuroanatomy &
  Allen Brain Atlas (mouse/rat gene expression) \\
\midrule
Disease ontology &
  Human Phenotype Ontology (HPO), Monarch Initiative \\
\midrule
Pathway analysis &
  Reactome, KEGG, WikiPathways, Enrichr \\
\midrule
Gene expression &
  GTEx, Human Protein Atlas (HPA) \\
\midrule
Epigenomics &
  GEO, ENCODE \\
\end{longtable}
}

\subsection*{SI-1.4\quad Query Design Rationale}

\paragraph{Broad initial queries.}
The first query to each literature database used broad terms (e.g.,
\texttt{"emotional dysregulation transdiagnostic neurobiology review"}) to
establish baseline coverage and identify landmark reviews.  Overly specific
Boolean conjunctions (e.g., combining four MeSH terms with AND) risk
returning zero results because different databases tokenize and index
differently; ToolUniverse wraps heterogeneous APIs, so robust recall
requires tolerant query design.

\paragraph{Domain-specific follow-ups.}
After inspecting Wave~1 returns, targeted follow-up queries were constructed
to fill gaps:
\begin{itemize}[nosep]
  \item Gross process model and DERS validation (measurement constructs).
  \item ADHD executive function overlap (transdiagnostic breadth).
  \item Linehan biosocial theory (theoretical foundations).
  \item Oxytocin and cortisol pathways (neuroendocrine mechanisms).
  \item FKBP5 chaperone cycle (stress biology).
\end{itemize}

\paragraph{Avoiding over-specified AND queries.}
Because ToolUniverse's literature tools pass queries as free-text strings
(not structured Boolean), adding excessive conjuncts can cause the
underlying API to return zero hits when any single term fails to match an
index entry.  We therefore preferred short, focused queries (2--4 key
concepts) and relied on result-count feedback to calibrate specificity.


% ============================================================================
% SI-2: COMPLETE TOOL CALL LOG (ALL 54 CALLS)
% ============================================================================
\sectitle{SI-2: Complete Tool Call Log}

The following table records every tool call in chronological order across
all five waves of analysis.  ``N'' reports the number of records returned
by the API (--- indicates an error or infrastructure failure).  ``Key
Finding'' summarizes the most salient result.  Arguments are reproduced
as abbreviated JSON; strings exceeding 60 characters are truncated with
\texttt{[...]}.

{\small
\setlength{\LTleft}{-1cm}
\setlength{\LTright}{-1cm}
\begin{longtable}{@{}c p{2.2cm} p{4.5cm} c p{4cm} c@{}}
\toprule
\textbf{\#} & \textbf{Tool} & \textbf{Arguments (JSON)} &
\textbf{N} & \textbf{Key Finding} & \textbf{W} \\
\midrule
\endfirsthead
\toprule
\textbf{\#} & \textbf{Tool} & \textbf{Arguments (JSON)} &
\textbf{N} & \textbf{Key Finding} & \textbf{W} \\
\midrule
\endhead
\bottomrule
\endlastfoot

% ---- Wave 1 (17 calls) ----

1 &
\texttt{PubMed\_search\_\allowbreak articles} &
\texttt{\{"query": "emotional dysregulation transdiag.\ [...] review", "max\_results": 10\}} &
1 &
Single review; confirmed need for broader queries &
1 \\
\midrule

2 &
\texttt{PubMed\_search\_\allowbreak articles} &
\texttt{\{"query": "emotion regulation prefrontal amygdala fMRI", "max\_results": 10\}} &
20 &
Silverman et al.\ 2019: vmPFC--amygdala connectivity &
1 \\
\midrule

3 &
\texttt{PubMed\_search\_\allowbreak articles} &
\texttt{\{"query": "DBT emotional dysregulation RCT", "max\_results": 10\}} &
20 &
Goldstein 2024 (\emph{JAMA Psychiatry}): DBT-A vs.\ TAU &
1 \\
\midrule

4 &
\texttt{openalex\_\allowbreak literature\_\allowbreak search} &
\texttt{\{"search\_keywords": "emotional dysregulation [...] neurocircuitry", "max\_results": 10, "year\_from": 2018\}} &
10 &
Beauchaine \& Cicchetti 2019: transdiagnostic framework &
1 \\
\midrule

5 &
\texttt{EuropePMC\_\allowbreak search\_\allowbreak articles} &
\texttt{\{"query": "emotional dysregulation amygdala [...] neuroimaging", "limit": 10\}} &
10 &
Etkin et al.\ 2015: common neural substrates &
1 \\
\midrule

6 &
\texttt{Semantic\allowbreak Scholar\_\allowbreak search\_papers} &
\texttt{\{"query": "emotional dysregulation transdiag.\ psychopathology", "limit": 10\}} &
10 &
Sloan et al.\ 2017: transdiagnostic $k$-factor model &
1 \\
\midrule

7 &
\texttt{Crossref\_\allowbreak search\_works} &
\texttt{\{"query": "emotional dysregulation [...] neurobiology", "limit": 10, "filter": "type:journal-article"\}} &
10 &
Aldao et al.\ 2010 meta-analysis (216 studies) &
1 \\
\midrule

8 &
\texttt{gwas\_search\_\allowbreak studies} &
\texttt{\{"disease\_trait": "neuroticism", "size": 10\}} &
10 &
GCST006476 ($N{=}449{,}484$): 136 loci &
1 \\
\midrule

9 &
\texttt{gwas\_search\_\allowbreak associations} &
\texttt{\{"disease\_trait": "neuroticism", "size": 10\}} &
--- &
EFO mapping error: returned ulcerative colitis (see SI-5) &
1 \\
\midrule

10 &
\texttt{ClinicalTrials\_\allowbreak search\_studies} &
\texttt{\{"query": "emotional dysregulation DBT treatment", "max\_results": 10\}} &
84 &
84 registered trials; majority phase II/III &
1 \\
\midrule

11 &
\texttt{UniProt\_get\_\allowbreak function\_by\_\allowbreak accession} &
\texttt{\{"accession": "P31645"\}} &
1 &
SLC6A4: Na$^+$/Cl$^-$-dependent serotonin reuptake &
1 \\
\midrule

12 &
\texttt{AllenBrain\_\allowbreak search\_\allowbreak structures} &
\texttt{\{"name": "amygdala", "num\_rows": 10\}} &
88 &
88 amygdalar sub-regions; basolateral complex dominant &
1 \\
\midrule

13 &
\texttt{AllenBrain\_\allowbreak search\_genes} &
\texttt{\{"gene\_acronym": "Slc6a4"\}} &
3 &
Gene ID 15342; mouse and rat expression data &
1 \\
\midrule

14 &
\texttt{kegg\_search\_\allowbreak pathway} &
\texttt{\{"keyword": "serotonin"\}} &
1 &
map07211: Serotonergic synapse &
1 \\
\midrule

15 &
\texttt{get\_HPO\_ID\_\allowbreak by\_phenotype} &
\texttt{\{"query": "emotional dysregulation", "limit": 5\}} &
2 &
HP:0100851 (Abnormal emotion); HP:0031467 &
1 \\
\midrule

16 &
\texttt{FAERS\_count\_\allowbreak reactions\_by\_\allowbreak drug\_event} &
\texttt{\{"medicinalproduct": "fluoxetine"\}} &
0 &
Empty count; likely API format issue (see SI-5) &
1 \\
\midrule

17 &
\texttt{PubChem\_get\_\allowbreak CID\_by\_\allowbreak compound\_name} &
\texttt{\{"name": "fluoxetine"\}} &
1 &
CID 3386 confirmed for fluoxetine &
1 \\
\midrule

% ---- Wave 2 (5 calls) ----

18 &
\texttt{kegg\_search\_\allowbreak pathway} &
\texttt{\{"keyword": "oxytocin"\}} &
1 &
map04921: Oxytocin signaling pathway &
2 \\
\midrule

19 &
\texttt{STRING\_\allowbreak functional\_\allowbreak enrichment} &
\texttt{\{"protein\_ids": ["SLC6A4", "FKBP5", [...]], "species": 9606, "category": "Process"\}} &
15 &
GO:0050890 ``cognition'' (5/8 genes, FDR\,=\,0.011) &
2 \\
\midrule

20 &
\texttt{Reactome\_map\_\allowbreak uniprot\_to\_\allowbreak pathways} &
\texttt{\{"uniprot\_id": "P31645"\}} &
2 &
R-HSA-380615, R-HSA-442660 (neurotransmitter release) &
2 \\
\midrule

21 &
\texttt{PubChem\_get\_\allowbreak compound\_\allowbreak properties\_\allowbreak by\_CID} &
\texttt{\{"cid": "3386", "properties": "MolecularFormula, [...] SMILES"\}} &
1 &
C$_{17}$H$_{18}$F$_3$NO; MW\,=\,309.33\,g/mol &
2 \\
\midrule

22 &
\texttt{kegg\_search\_\allowbreak pathway} &
\texttt{\{"query": "serotonin"\}} $\to$ retried as \texttt{\{"keyword": "serotonin"\}} &
1 &
Validation caught wrong param name; retry succeeded &
2 \\
\midrule

% ---- Wave 3 (10 calls) ----

23 &
\texttt{PubMed\_search\_\allowbreak articles} &
\texttt{\{"query": "Gross process model emotion regulation", "max\_results": 5\}} &
5 &
Gross 2002 (process model); John \& Gross 2004 &
3 \\
\midrule

24 &
\texttt{PubMed\_search\_\allowbreak articles} &
\texttt{\{"query": "DERS validation Difficulties Emotion Regulation Scale", "max\_results": 10\}} &
11 &
Gratz \& Roemer 2004 (original DERS) &
3 \\
\midrule

25 &
\texttt{PubMed\_search\_\allowbreak articles} &
\texttt{\{"query": "ADHD executive function emotional dysregulation", "max\_results": 5\}} &
4 &
Pozo-Rodriguez et al.\ 2026: EF--ED comorbidity &
3 \\
\midrule

26 &
\texttt{gwas\_search\_\allowbreak studies} &
\texttt{\{"disease\_trait": "anxiety", "size": 5\}} &
11 &
GCST90319327 ($N{=}486{,}919$) &
3 \\
\midrule

27 &
\texttt{gwas\_search\_\allowbreak studies} &
\texttt{\{"disease\_trait": "depression", "size": 5\}} &
27 &
27 GWAS; largest $N{=}1.35$\,M; UKB overlap &
3 \\
\midrule

28 &
\texttt{kegg\_search\_\allowbreak pathway} &
\texttt{\{"keyword": "cortisol"\}} &
1 &
map04927: Cortisol synthesis and secretion &
3 \\
\midrule

29 &
\texttt{Reactome\_map\_\allowbreak uniprot\_to\_\allowbreak pathways} &
\texttt{\{"uniprot\_id": "Q13451"\}} &
1 &
R-HSA-3371497: HSP90 chaperone cycle (FKBP5) &
3 \\
\midrule

30 &
\texttt{STRING\_\allowbreak functional\_\allowbreak enrichment} &
\texttt{\{"protein\_ids": ["SLC6A4", "FKBP5", [...]], "species": 9606, "category": "KEGG"\}} &
4 &
Alcoholism (FDR\,=\,0.008); Serotonergic synapse (0.012) &
3 \\
\midrule

31 &
\texttt{openalex\_\allowbreak literature\_\allowbreak search} &
\texttt{\{"search\_keywords": "emotion regulation neuroscience", "max\_results": 5\}} &
5 &
Troy et al.\ 2022 (flexibility); Kredlow et al.\ 2021 &
3 \\
\midrule

32 &
\texttt{PubMed\_search\_\allowbreak articles} &
\texttt{\{"query": "Linehan biosocial theory emotional dysreg.\ [...] borderline", "max\_results": 10\}} &
10 &
Crowell et al.\ 2009: developmental biosocial model &
3 \\
\midrule

% ---- Wave 4 (14 calls) ----

33 &
\texttt{PubMed MCP\_\allowbreak search\_\allowbreak articles} &
\texttt{\{"query": "emotional dysregulation transdiag.\ meta-analysis"\}} &
--- &
Proxy error; fell back to ToolUniverse (Call~52) &
4 \\
\midrule

34 &
\texttt{PubMed MCP\_\allowbreak search\_\allowbreak articles} &
\texttt{\{"query": "emotion regulation sex differences [...] fMRI"\}} &
--- &
Proxy error; Anthropic MCP proxy unreachable &
4 \\
\midrule

35 &
\texttt{PubMed MCP\_\allowbreak search\_\allowbreak articles} &
\texttt{\{"query": "SLC6A4 FKBP5 BDNF methylation [...] dysregulation"\}} &
--- &
Proxy error; re-issued via ToolUniverse (Call~53) &
4 \\
\midrule

36 &
\texttt{PubMed MCP\_\allowbreak find\_related\_\allowbreak articles} &
\texttt{\{"pmids": [12212647, 15509284, 19379027, 29942085]\}} &
--- &
Proxy error; related-article traversal abandoned &
4 \\
\midrule

37 &
\texttt{bioRxiv MCP\_\allowbreak search\_\allowbreak preprints} &
\texttt{\{"subject": "neuroscience", "days\_back": 90\}} &
30 &
General neuroscience preprints; no ED-specific results &
4 \\
\midrule

38 &
\texttt{medRxiv MCP\_\allowbreak search\_\allowbreak preprints} &
\texttt{\{"category": "psychiatry and clinical psychology"\}} &
--- &
Validation error: category not in supported enum &
4 \\
\midrule

39 &
\texttt{GTEx\_get\_\allowbreak median\_gene\_\allowbreak expression} &
\texttt{\{"gencode\_id": "ENSG00000108576.14"\}} &
0 &
SLC6A4 not expressed in standard GTEx brain tissues &
4 \\
\midrule

40 &
\texttt{DGIdb\_get\_\allowbreak drug\_gene\_\allowbreak interactions} &
\texttt{\{"genes": ["SLC6A4", "FKBP5", "COMT", "MAOA", "OXTR", "CRH", "NR3C1", "BDNF"]\}} &
464 &
464 drug--gene interactions; SSRIs dominate SLC6A4 &
4 \\
\midrule

41 &
\texttt{Enrichr\_GO\_\allowbreak Biological\_\allowbreak Process\_2023} &
\texttt{\{"genes": ["SLC6A4", "FKBP5", [...] "BDNF"]\}} &
--- &
Top: response to stress (111K chars, saved to file) &
4 \\
\midrule

42 &
\texttt{WikiPathways\_\allowbreak find\_pathways\_\allowbreak by\_gene} &
\texttt{\{"gene": "SLC6A4"\}} &
10 &
WP5420 (ADHD/ASD); WP3944 (serotonin receptor 4/6/7) &
4 \\
\midrule

43 &
\texttt{Monarch\_\allowbreak search\_gene} &
\texttt{\{"query": "SLC6A4"\}} &
1 &
HGNC:11050, ENSG00000108576, OMIM:182138 &
4 \\
\midrule

44 &
\texttt{PharmGKB\_\allowbreak search\_genes} &
\texttt{\{"query": "SLC6A4"\}} &
1 &
PA312; chr17q11.2; extensive pharmacogenomic annotation &
4 \\
\midrule

45 &
\texttt{HPA\_get\_rna\_\allowbreak expression\_\allowbreak by\_source} &
\texttt{\{"gene\_name": "SLC6A4", "tissue": "brain", "region": "amygdala"\}} &
0 &
No expression in amygdala (informative negative) &
4 \\
\midrule

46 &
\texttt{GEO\_search\_\allowbreak methylation\_\allowbreak datasets} &
\texttt{\{"query": "FKBP5 methylation PTSD"\}} &
0 &
Query too specific; no matching datasets &
4 \\
\midrule

% ---- Wave 5 (8 calls) ----

47 &
\texttt{GTEx\_get\_\allowbreak median\_gene\_\allowbreak expression} &
\texttt{\{"operation": "median\_gene\_expression", "gencode\_id": "ENSG00000108576.14"\}} &
0 &
Confirmed: SLC6A4 absent from GTEx brain panel &
5 \\
\midrule

48 &
\texttt{WikiPathways\_\allowbreak find\_pathways\_\allowbreak by\_gene} &
\texttt{\{"gene": "SLC6A4"\}} &
10 &
Confirmed 10 pathways; WP5420 validated &
5 \\
\midrule

49 &
\texttt{HPA\_get\_rna\_\allowbreak expression\_\allowbreak by\_source} &
\texttt{\{"gene\_name": "SLC6A4", "tissue": "brain", "region": "amygdala"\}} &
0 &
Confirmed: no HPA data for SLC6A4 in amygdala &
5 \\
\midrule

50 &
\texttt{Monarch\_get\_\allowbreak gene\_diseases} &
\texttt{\{"subject": "HGNC:11050"\}} &
--- &
Parameter error: \texttt{subject} not accepted &
5 \\
\midrule

51 &
\texttt{DGIdb\_get\_\allowbreak gene\_\allowbreak druggability} &
\texttt{\{"genes": ["SLC6A4", "FKBP5", "COMT", "MAOA", "OXTR", "CRH", "NR3C1", "BDNF"]\}} &
8 &
All 8 genes in DRUGGABLE GENOME &
5 \\
\midrule

52 &
\texttt{PubMed\_search\_\allowbreak articles} &
\texttt{\{"query": "emotional dysregulation transdiag.\ meta-analysis", "max\_results": 30\}} &
14 &
14 meta-analyses; Aldao 2010 confirmed; Sloan 2017 &
5 \\
\midrule

53 &
\texttt{PubMed\_search\_\allowbreak articles} &
\texttt{\{"query": "SLC6A4 polymorphism SSRI treatment response", "max\_results": 20\}} &
19 &
5-HTTLPR moderates SSRI efficacy &
5 \\
\midrule

54 &
\texttt{PubMed\_search\_\allowbreak articles} &
\texttt{\{"query": "emotion regulation sex differences brain neuroimaging", "max\_results": 20\}} &
20 &
Females: greater amygdala--PFC connectivity during reappraisal &
5 \\

\end{longtable}
}

\noindent
\textbf{Summary statistics.}
54 tool calls across 5 waves.  23 unique tool types.  13 databases queried.
6 infrastructure failures (proxy errors, validation errors, schema mismatches).
3 informative negatives.  3 queries saved to file due to output exceeding
50{,}000 characters (DGIdb interactions, Enrichr enrichment, DGIdb
druggability).  42 calls (78\%) returned usable results contributing
directly to the main text.


% ============================================================================
% SI-3: STATISTICAL METHODS
% ============================================================================
\sectitle{SI-3: Statistical Methods}

\subsection*{SI-3.1\quad STRING Functional Enrichment}

Protein--protein interaction network enrichment was performed using STRING
v12 (Szklarczyk et al., \textit{Nucleic Acids Res.}, 2023) via the ToolUniverse wrapper
\texttt{STRING\_functional\_enrichment}.

\paragraph{Input gene set.}
Eight genes were selected based on convergent evidence from the literature
search (Wave~1): \texttt{SLC6A4}, \texttt{FKBP5}, \texttt{COMT},
\texttt{MAOA}, \texttt{OXTR}, \texttt{CRH}, \texttt{NR3C1}, and
\texttt{BDNF}.  These were chosen because each appeared in at least three
independent database returns (PubMed, GWAS Catalog, and UniProt/Reactome)
in the context of emotional regulation, stress response, or affective
disorder neurobiology.

\paragraph{Statistical test.}
STRING computes enrichment using a hypergeometric test.  For a gene set of
size $k = 8$ drawn from the background of all human protein-coding genes
in STRING ($N \approx 19{,}566$), the probability that $m$ or more genes
fall into a given GO category of size $M$ is:

\begin{equation}
  p = 1 - \sum_{i=0}^{m-1} \frac{\binom{M}{i}\binom{N-M}{k-i}}{\binom{N}{k}}
\end{equation}

\paragraph{Multiple testing correction.}
The Benjamini--Hochberg procedure was applied to control the false discovery
rate (FDR) across all tested categories within each enrichment run.

\begin{itemize}[nosep]
  \item \textbf{GO Biological Process} (Call~19): 15 enriched terms returned;
        all with FDR\,$<$\,0.05.  Top term: GO:0050890 ``cognition'' (5 of
        8 input genes, FDR\,=\,0.011).
  \item \textbf{KEGG Pathways} (Call~30): 4 enriched pathways returned; all
        with FDR\,$<$\,0.05.  Top pathway: ``Alcoholism'' (FDR\,=\,0.008),
        followed by ``Serotonergic synapse'' (FDR\,=\,0.012).
\end{itemize}

\subsection*{SI-3.2\quad Enrichr Cross-Validation}

Independent enrichment analysis was performed via the Enrichr API (Call~41)
using the Fisher exact test.  The same 8-gene input set was submitted to
the GO Biological Process 2023 library.  Enrichr's enrichment pipeline
serves as a cross-validation of the STRING results: both tools use different
background sets and statistical implementations, yet converge on
stress-response and behavioral GO categories.

\subsection*{SI-3.3\quad GWAS Summary Statistics}

GWAS studies were retrieved from the NHGRI-EBI GWAS Catalog.  We did not
perform novel statistical analyses on GWAS summary statistics; we report
published effect sizes, $p$-values, and sample sizes as catalogued.  Cross-trait
comparisons (neuroticism, depression, anxiety) rely on the observation that the
same UK Biobank (UKB) cohort contributes to multiple studies, enabling
qualitative assessment of genetic overlap.

\subsection*{SI-3.4\quad DGIdb Drug--Gene Interactions}

DGIdb (Call~40) returns curated drug--gene interaction counts aggregated
from published literature and clinical databases.  No statistical test
is applied; the 464 interactions are descriptive.  Interaction counts
include all evidence levels and are not filtered by clinical validation
stage.

\subsection*{SI-3.5\quad Limitations of the Statistical Approach}

\begin{enumerate}[nosep]
  \item \textbf{Confirmation bias in gene selection.}
        The 8-gene set was assembled from prior literature, not from an
        unbiased genome-wide screen.  Enrichment of these genes in
        behavior-related GO categories is therefore partly tautological:
        genes selected for their roles in affective neurobiology are expected
        to enrich for behavioral processes.  This enrichment should be
        interpreted as a consistency check, not as novel discovery.
  \item \textbf{Background set sensitivity.}
        The significance of hypergeometric enrichment depends on the size
        and composition of the background set.  STRING's default background
        ($\sim$19,566 genes) includes all annotated human proteins.  A smaller,
        brain-specific background could yield different FDR values.
  \item \textbf{Multiple comparisons across waves.}
        Across 5 waves, 15 GO + 4 KEGG terms were tested; all pass
        FDR\,$<$\,0.05 within their respective databases.  We did not apply
        formal correction across waves; instead, we require that all
        reported enrichments independently pass FDR\,$<$\,0.05.
  \item \textbf{No novel meta-analysis.}
        This review does not perform a quantitative meta-analysis of effect
        sizes.  All statistics reported are sourced directly from the
        original publications or database APIs.
\end{enumerate}


% ============================================================================
% SI-4: EVIDENCE GRADING
% ============================================================================
\sectitle{SI-4: Evidence Grading}

Each major claim in the main text is graded below according to the following
schema:

\begin{itemize}[nosep]
  \item \textbf{Strong}: Supported by $\geq$2 RCTs, a meta-analysis, or a
        large-cohort GWAS ($N > 100{,}000$).
  \item \textbf{Moderate}: Supported by a single large study ($N > 500$)
        or consistent observational evidence with cross-validation.
  \item \textbf{Preliminary}: Based on small-$N$ studies, preprints, or
        indirect inference.
\end{itemize}

\subsection*{Claim 1: ``Neuroticism reflects regulation failure, not
generation excess''}

\begin{description}[style=nextline, nosep]
  \item[Grade] \textbf{Strong.}
  \item[Evidence chain]
    Silverman et al.\ 2019 ($N > 500$), retrieved via PubMed (Call~2) and
    cross-verified in OpenAlex (Call~31).  Replicated by Meiering et al.\
    2026 large-sample neuroimaging.  Neuroticism correlated with reduced
    vmPFC--amygdala functional connectivity during reappraisal, with no
    significant difference in amygdala reactivity.
\end{description}

\subsection*{Claim 2: ``Eight candidate genes enrich for behavioral and
cognitive GO processes''}

\begin{description}[style=nextline, nosep]
  \item[Grade] \textbf{Moderate.}
  \item[Evidence chain]
    STRING v12 functional enrichment (Call~19): 5 of 8 genes annotated to
    GO:0050890 ``cognition'' (FDR\,=\,0.011).  Cross-validated by Enrichr
    (Call~41): independent Fisher exact test converged on stress-response
    GO terms.  However, the gene set is hypothesis-driven (see SI-3.5),
    limiting the strength of the inference.
\end{description}

\subsection*{Claim 3: ``DBT is effective across diagnostic boundaries''}

\begin{description}[style=nextline, nosep]
  \item[Grade] \textbf{Strong.}
  \item[Evidence chain]
    Three independent RCTs from PubMed (Call~3) and ClinicalTrials.gov
    (Call~10):
    (i)~Goldstein et al.\ 2024 (\emph{JAMA Psychiatry}): DBT-A for
    adolescent bipolar disorder;
    (ii)~Bemmouna et al.\ 2025: DBT for autistic adults;
    (iii)~Norman-Nott et al.\ 2025 (\emph{JAMA Network Open}): DBT for
    chronic pain.  Three RCTs across four diagnostic categories.
\end{description}

\subsection*{Claim 4: ``Polygenic architecture overlaps across affective
disorders''}

\begin{description}[style=nextline, nosep]
  \item[Grade] \textbf{Strong.}
  \item[Evidence chain]
    GWAS Catalog: neuroticism (Call~8, 20 studies), depression (Call~27,
    27 studies), anxiety (Call~26, 11 studies).  Key accessions:
    GCST006476 ($N{=}449{,}484$, 136 loci), GCST90029028
    ($N{=}1.35$\,M), GCST90319327 ($N{=}486{,}919$).  Same UKB cohort
    contributes to all three, with published $r_g > 0.5$.
\end{description}

\subsection*{Claim 5: ``SLC6A4 is not expressed in amygdala''}

\begin{description}[style=nextline, nosep]
  \item[Grade] \textbf{Moderate.}
  \item[Evidence chain]
    HPA (Calls~45, 49): no expression data for SLC6A4 in amygdala.
    GTEx (Calls~39, 47): SLC6A4 absent from standard brain panel.
    Allen Brain Atlas (Call~13): Slc6a4 expression concentrated in raphe
    nuclei.  This is an informative negative consistent with known
    neuroanatomy, but absence of evidence in expression atlases does not
    constitute evidence of absence at single-cell resolution.
\end{description}

\subsection*{Claim 6: ``464 drug--gene interactions across 8 candidate
genes''}

\begin{description}[style=nextline, nosep]
  \item[Grade] \textbf{Strong.}
  \item[Evidence chain]
    DGIdb (Call~40): directly queried 8 genes, returned 464 curated
    interactions.  All 8 genes confirmed in the DRUGGABLE GENOME
    (Call~51).  SSRIs dominate SLC6A4 interactions.
\end{description}

\subsection*{Claim 7: ``5-HTTLPR modulates SSRI response''}

\begin{description}[style=nextline, nosep]
  \item[Grade] \textbf{Moderate.}
  \item[Evidence chain]
    PubMed pharmacogenomics (Call~53): 19 papers on SLC6A4 polymorphism
    and SSRI treatment response, including Stein et al.\ 2021 meta-analysis.
    PharmGKB PA312 (Call~44) confirms extensive pharmacogenomic annotation.
    Effect sizes are small and contested in recent large-sample studies.
\end{description}


% ============================================================================
% SI-5: LIMITATIONS AND POTENTIAL BIASES
% ============================================================================
\sectitle{SI-5: Limitations and Potential Biases}

\subsection*{SI-5.1\quad Tool Availability Bias}

Only databases with ToolUniverse wrappers were queried.  Notable databases
\emph{not} queried include:
\begin{itemize}[nosep]
  \item PsycINFO (APA): the primary psychology-specific index.
  \item Cochrane Library: systematic reviews and meta-analyses.
  \item ENCODE / Roadmap Epigenomics: chromatin state data.
\end{itemize}

\subsection*{SI-5.2\quad Search Language Bias}

All queried databases are English-language or English-indexed.  Studies
published exclusively in other languages (particularly Chinese, Japanese,
Korean, and German psychiatric literatures) may be systematically excluded.

\subsection*{SI-5.3\quad LLM Planning Failure}

Tool\_Finder\_LLM returned an empty result set on both invocations due to
the 120-second Claude CLI timeout.  Fallback to \texttt{grep\_tools} and
\texttt{find\_tools} partially mitigates this, but may miss semantically
related tools with non-obvious names.  Post hoc review confirmed no
high-relevance tools were missed.

\subsection*{SI-5.4\quad GWAS EFO Mapping Error}

Call~9 returned ulcerative colitis associations instead of neuroticism due
to partial string matching in the EFO trait resolution.  Bug subsequently
fixed in the ToolUniverse codebase.  Erroneous data were identified on
inspection and discarded; correct data obtained from Call~8.

\subsection*{SI-5.5\quad FAERS Empty Results}

Call~16 returned count\,=\,0 for fluoxetine, inconsistent with known FAERS
data.  Likely cause: \texttt{medicinalproduct} field requires exact case
or salt name (``FLUOXETINE HYDROCHLORIDE'').  No adverse event claims made
in the main text.

\subsection*{SI-5.6\quad PubMed MCP Proxy Errors}

Calls 33--36 (all 4 PubMed MCP queries) failed due to Anthropic MCP proxy
being unreachable.  Three of four queries were recovered through ToolUniverse
PubMed endpoints (Calls~52--54).  Related-article traversal (Call~36) was
not retried but does not affect any claims.

\subsection*{SI-5.7\quad No Systematic Risk-of-Bias Assessment}

Individual RCTs were not subjected to Cochrane RoB-2 or equivalent
assessment.  This is a computationally augmented narrative review, not a
registered systematic review.

\subsection*{SI-5.8\quad Hallucination Controls}

\begin{enumerate}[nosep]
  \item \textbf{No unsourced claims.}  Every factual statement traces to a
        tool call in SI-2.
  \item \textbf{Cross-database verification.}  Key findings verified across
        $\geq$2 independent databases.
  \item \textbf{Verbatim argument logging.}  All tool call arguments recorded
        as JSON.
  \item \textbf{Anomaly documentation.}  All unexpected results documented
        with root cause analysis.
  \item \textbf{No model-knowledge-only claims.}  The review does not assert
        facts not returned by at least one database query.
\end{enumerate}

\subsection*{SI-5.9\quad Gene Selection Bias}

The 8-gene set was assembled from prior literature, not an unbiased screen.
STRING enrichment should be interpreted as a consistency check.  A truly
unbiased approach would begin with genome-wide significant loci from GWAS
studies and map them to genes---a natural extension for future work.

\subsection*{SI-5.10\quad SLC6A4 Expression Null}

GTEx (Calls~39, 47) and HPA (Calls~45, 49) returned no SLC6A4 expression
in brain tissues including amygdala.  This is informative but reflects atlas
limitations: GTEx does not dissect raphe nuclei, and HPA resolution may
miss low-abundance transcripts.  Absence of evidence $\neq$ evidence of
absence at single-cell resolution.

\subsection*{SI-5.11\quad DGIdb Interaction Counts}

The 464 drug--gene interactions (Call~40) include all evidence levels from
curated literature, not filtered by clinical validation stage.  Counts
should not be interpreted as 464 clinically validated drug targets.


% ============================================================================
% SI-6: AUTHOR CONTRIBUTIONS
% ============================================================================
\sectitle{SI-6: Author Contributions}

\subsection*{Chris Davis (Human Author)}

Conceived the review topic and research question.  Designed the multi-database
search strategy and iterative refinement protocol.  Directed the agentic
workflow, specifying which databases to query and in what order.  Provided
editorial guidance on narrative structure, figure design (rejecting two figure
iterations and directing the shift from circuit diagram to data-driven
enrichment plot), and prose style (specifying the Williams \emph{Style:
Lessons in Clarity and Grace} style guide).  Reviewed all outputs for factual
accuracy and coherence.  Approved the final manuscript.  C.D.\ is the sole
human author and takes responsibility for the intellectual content.

\subsection*{Claude Opus 4.6 (Anthropic, model \texttt{claude-opus-4-6}, 1M context)}

Executed all ToolUniverse database queries (54 tool calls across 5 waves).
Wrote the prose under C.D.'s editorial direction, applying the Williams style
guide.  Generated the Ti\textit{k}Z enrichment figure.  Constructed the
Bib\TeX{} bibliography from verified DOI/PMID data.  Compiled the \LaTeX{}
document.

\paragraph{Scope limitations.}
Claude Opus 4.6 did \emph{not}:
\begin{itemize}[nosep]
  \item independently select the research question;
  \item independently verify facts against primary sources outside the
        ToolUniverse pipeline;
  \item access any databases not documented in the tool call log (SI-2);
  \item exercise independent editorial judgment on which claims to include or
        exclude.
\end{itemize}

\subsection*{ToolUniverse (Harvard MIMS Lab, v1.1.5)}

Provided the standardized API layer through which all 1,700+ scientific
database queries were executed.  ToolUniverse did \emph{not} contribute to
study design, data interpretation, or manuscript writing.  It is an
open-source software tool, not a sentient contributor.  Its listing
acknowledges its essential role in enabling reproducible methodology.  If
journal policy prohibits non-human authorship, ToolUniverse should be cited
in the Methods section and Acknowledgments rather than the author byline.


% ============================================================================
% SI-7: PROMPTS AND AGENT INSTRUCTIONS
% ============================================================================
\sectitle{SI-7: Prompts and Agent Instructions}

Three writing agents were dispatched in parallel, each receiving a distinct
section assignment and a shared set of instructions.  All agents used the same
model (\texttt{claude-opus-4-6}, 1M token context window).  This section
reproduces the key directives from each prompt to ensure transparency and
reproducibility.

\subsection*{SI-7.1\quad Agent 1: Introduction + Neurocircuit Model}

\begin{description}[style=nextline, nosep]
  \item[Assignment]
    Write expanded \LaTeX{} sections for an academic review paper on emotional
    dysregulation.
  \item[Key directive]
    ``Apply Williams \emph{Style: Lessons in Clarity and Grace} principles.
    Use active voice.  Keep sentences under 25 words where possible.  Vary
    sentence length for rhythm.  Place stress position (the important idea) at
    the end of each sentence.''
  \item[Evidence provided]
    \begin{itemize}[nosep]
      \item Amygdala--PFC connectivity findings (Silverman et al.\ 2019)
      \item Gross process model (Gross 2002; John \& Gross 2004)
      \item Transdiagnostic framework (Beauchaine \& Cicchetti 2019)
      \item Etkin et al.\ 2015: common neural substrates
      \item DERS validation (Gratz \& Roemer 2004)
      \item Linehan biosocial theory (Crowell et al.\ 2009)
    \end{itemize}
  \item[Approximate token budget] 3,000--4,000 tokens output
  \item[Model] \texttt{claude-opus-4-6} (1M context)
\end{description}

\subsection*{SI-7.2\quad Agent 2: Genetic Architecture + Therapeutic Landscape}

\begin{description}[style=nextline, nosep]
  \item[Assignment]
    Write expanded \LaTeX{} sections covering genetic architecture, epigenetic
    mechanisms, pharmacogenomics, and the therapeutic landscape.
  \item[Key directive]
    ``Integrate GWAS, STRING enrichment, and Reactome pathway data into a
    coherent narrative.  Use the enrichment figure as an anchor.  Do not
    overclaim: distinguish between genome-wide evidence and candidate-gene
    consistency checks.''
  \item[Evidence provided]
    \begin{itemize}[nosep]
      \item GWAS Catalog data: neuroticism (136 loci), depression (27 studies),
            anxiety (11 studies)
      \item STRING enrichment: GO:0050890 cognition (FDR\,=\,0.011)
      \item Reactome: SLC6A4 transporter pathways, FKBP5 chaperone cycle
      \item WikiPathways WP5420 (ADHD/ASD pathways)
      \item DGIdb: 464 drug--gene interactions, druggable genome annotations
      \item PharmGKB PA312 (SLC6A4 pharmacogenomics)
      \item DBT RCTs: Goldstein 2024, Bemmouna 2025, Norman-Nott 2025
      \item ClinicalTrials.gov: 84 registered trials
      \item PubMed pharmacogenomics: 19 papers on 5-HTTLPR and SSRI response
    \end{itemize}
  \item[Approximate token budget] 3,000--4,000 tokens output
  \item[Model] \texttt{claude-opus-4-6} (1M context)
\end{description}

\subsection*{SI-7.3\quad Agent 3: Supplementary Information}

\begin{description}[style=nextline, nosep]
  \item[Assignment]
    Write the complete Supplementary Information appendix for all waves of
    analysis.
  \item[Key directive]
    ``Document every tool call with verbatim JSON arguments.  Be exhaustive
    rather than selective.  All negative and null results must be recorded
    with root cause analysis.  This appendix must be sufficient for an
    independent researcher to reproduce the entire analysis.''
  \item[Structure specified]
    Six sections: (1)~Search Strategy, (2)~Tool Call Log, (3)~Statistical
    Methods, (4)~Evidence Grading, (5)~Limitations and Potential Biases,
    (6)~Data Availability.  Expanded in the second harvest to include
    SI-7 through SI-9.
  \item[Approximate token budget] 5,000--8,000 tokens output
  \item[Model] \texttt{claude-opus-4-6} (1M context)
\end{description}

\subsection*{SI-7.4\quad Parallel Execution}

All three agents ran in parallel, receiving their prompts simultaneously.
Each agent operated on a non-overlapping section of the manuscript.  No agent
had access to the output of any other agent during generation.  Post-hoc
editorial integration (deduplication, cross-reference alignment, and style
harmonization) was performed by C.D.


% ============================================================================
% SI-8: NEGATIVE AND NULL RESULTS
% ============================================================================
\sectitle{SI-8: Negative and Null Results}

Transparent documentation of negative and null results guards against
survivorship bias and ensures reproducibility.  This section catalogs every
query that returned empty, unexpected, or erroneous results.

\subsection*{SI-8.1\quad Infrastructure Failures}

Infrastructure failures reflect problems with tool availability, network
connectivity, or parameter validation---not with the underlying scientific
databases.

{\small
\begin{longtable}{@{}c p{2.8cm} p{2.8cm} p{5.8cm}@{}}
\toprule
\textbf{Call} & \textbf{Tool} & \textbf{Failure Mode} &
\textbf{Root Cause and Impact} \\
\midrule
\endfirsthead
\toprule
\textbf{Call} & \textbf{Tool} & \textbf{Failure Mode} &
\textbf{Root Cause and Impact} \\
\midrule
\endhead
\bottomrule
\endlastfoot

--- &
\texttt{Tool\_Finder\_LLM} (both invocations) &
Empty result set &
Claude CLI 120\,s timeout over 1,700-tool catalogue.  \textbf{Impact}: Fell
back to \texttt{grep\_tools}/\texttt{find\_tools}; post-hoc review confirmed
no tools missed. \\
\midrule

--- &
\texttt{BioRxiv\_search\_\allowbreak preprints} (TU) &
Tool does not exist &
Ghost reference: no keyword-search endpoint in ToolUniverse for bioRxiv.
MCP-native bioRxiv (Call~37) used instead.  \textbf{Impact}: None. \\
\midrule

9 &
\texttt{gwas\_search\_\allowbreak associations} &
Wrong trait returned &
EFO partial string matching mapped ``neuroticism'' to ulcerative colitis.
Bug fixed in ToolUniverse.  \textbf{Impact}: Discarded; correct data from
Call~8. \\
\midrule

--- &
\texttt{PubMed\_search\_\allowbreak articles} &
0 results &
Over-specified 7-term AND conjunction.  \textbf{Impact}: Decomposed into
shorter queries in Waves~3--5. \\
\midrule

16 &
\texttt{FAERS\_count\_\allowbreak reactions\_by\_\allowbreak drug\_event} &
count\,=\,0 &
Likely query format issue (case or salt name).  \textbf{Impact}: No FAERS
claims in main text. \\
\midrule

--- &
\texttt{proteins\_api\_search} &
HTTP 400 &
Endpoint requires structured parameters.  \textbf{Impact}: Data from UniProt
(Call~11) and STRING (Call~19). \\
\midrule

--- &
\texttt{PubMed\_Guidelines\_\allowbreak Search} &
0 results &
Relevance filter bug; fixed in ToolUniverse.  \textbf{Impact}: Guidelines
from standard PubMed queries. \\
\midrule

33--36 &
\texttt{PubMed MCP} (4 queries) &
Proxy error &
Anthropic MCP proxy unreachable.  \textbf{Impact}: 3 of 4 recovered via
ToolUniverse (Calls~52--54). \\
\midrule

38 &
\texttt{medRxiv MCP} &
Validation error &
Category not in supported enumeration.  \textbf{Impact}: medRxiv not
systematically searched (coverage gap). \\
\midrule

50 &
\texttt{Monarch\_get\_\allowbreak gene\_diseases} &
Parameter error &
\texttt{subject} not accepted; schema mismatch.  \textbf{Impact}: Gene
metadata from \texttt{Monarch\_search\_gene} (Call~43). \\

\end{longtable}
}

\subsection*{SI-8.2\quad Informative Negatives}

Informative negatives are scientifically meaningful null results that
constrain biological interpretation.

{\small
\begin{longtable}{@{}c p{2.8cm} p{2.8cm} p{5.8cm}@{}}
\toprule
\textbf{Call} & \textbf{Tool} & \textbf{Null Result} &
\textbf{Biological Interpretation} \\
\midrule
\endfirsthead
\toprule
\textbf{Call} & \textbf{Tool} & \textbf{Null Result} &
\textbf{Biological Interpretation} \\
\midrule
\endhead
\bottomrule
\endlastfoot

39, 47 &
\texttt{GTEx\_get\_median\_\allowbreak gene\_expression} &
SLC6A4 not expressed in standard GTEx brain regions &
SLC6A4 expression is concentrated in the dorsal and median raphe nuclei,
which are not among the standard GTEx brain dissection regions.  Consistent
with serotonergic neuroanatomy: transporter expression restricted to
raphe cell bodies, not projection targets.  \textbf{Impact}: Strengthens
raphe-specific localization claims. \\
\midrule

45, 49 &
\texttt{HPA\_get\_rna\_\allowbreak expression\_\allowbreak by\_source} &
No SLC6A4 expression in amygdala &
Corroborates GTEx finding.  Consistent with Allen Brain Atlas mouse data
(Call~13) showing Slc6a4 in raphe nuclei.  \textbf{Impact}: No main-text
claims assert amygdalar SLC6A4 expression. \\
\midrule

16 &
\texttt{FAERS\_count\_\allowbreak reactions\_by\_\allowbreak drug\_event} &
0 adverse events for fluoxetine &
Query format issue, not biological null.  \textbf{Impact}: No FAERS-based
safety claims. \\
\midrule

46 &
\texttt{GEO\_search\_\allowbreak methylation\_\allowbreak datasets} &
0 results for ``FKBP5 methylation PTSD'' &
Query too specific; GEO datasets annotated with broader keywords.  FKBP5
epigenetic evidence drawn from PubMed literature instead.
\textbf{Impact}: No direct epigenomic dataset analysis in main text. \\

\end{longtable}
}

\subsection*{SI-8.3\quad Summary Assessment}

Of the 54 total tool calls:

\begin{itemize}[nosep]
  \item \textbf{42 calls} (78\%) returned usable results contributing
        directly to the main text or enrichment analysis.
  \item \textbf{6 calls} (11\%) were infrastructure failures with no impact
        on conclusions (data recovered through alternative endpoints).
  \item \textbf{3 calls} (6\%) returned informative negatives that
        strengthened biological interpretation (SLC6A4 raphe-specific
        expression).
  \item \textbf{3 calls} (6\%) returned zero results due to overly specific
        queries or API format issues; targeted data obtained through
        refined follow-up queries.
\end{itemize}

\noindent
No claim in the main text relies solely on data from a failed or null query.
Every factual assertion is supported by at least one successful database
return, and key claims are cross-verified across two or more independent
databases (see SI-4).


% ============================================================================
% SI-9: DATA AVAILABILITY
% ============================================================================
\sectitle{SI-9: Data Availability}

All queries documented in SI-2 are reproducible via the ToolUniverse MCP
protocol (v1.1.5).  The following identifiers and accession numbers are
provided to enable independent replication.

\subsection*{SI-9.1\quad Software and Protocol}

\begin{itemize}[nosep]
  \item \textbf{ToolUniverse}: v1.1.5,
        \url{https://github.com/mims-harvard/ToolUniverse}
  \item \textbf{Protocol}: MCP (Model Context Protocol) stdio mode
  \item \textbf{Configuration}: \texttt{TOOLUNIVERSE\_CACHE\_ENABLED=true},
        \texttt{TOOLUNIVERSE\_COERCE\_TYPES=true},
        \texttt{TOOLUNIVERSE\_STRICT\_VALIDATION=false}
\end{itemize}

\subsection*{SI-9.2\quad STRING Enrichment Input}

\begin{verbatim}
{
  "protein_ids": ["SLC6A4","FKBP5","COMT","MAOA",
                   "OXTR","CRH","NR3C1","BDNF"],
  "species": 9606,
  "category": "Process"
}
\end{verbatim}

\noindent
For KEGG enrichment, the same protein list was used with
\texttt{"category": "KEGG"}.

\subsection*{SI-9.3\quad GWAS Catalog Accession IDs}

{\small
\begin{longtable}{@{}l l r l@{}}
\toprule
\textbf{Accession} & \textbf{Trait} & \textbf{Sample Size} & \textbf{Reference} \\
\midrule
\endfirsthead
\toprule
\textbf{Accession} & \textbf{Trait} & \textbf{Sample Size} & \textbf{Reference} \\
\midrule
\endhead
\bottomrule
\endlastfoot
GCST006476  & Neuroticism & 449,484  & Nagel et al.\ 2018 \\
\midrule
GCST007339  & Neuroticism & 523,783  & Baselmans et al.\ 2019 \\
\midrule
GCST90029028 & Depression & 1,350,000 & Howard et al.\ 2019 \\
\midrule
GCST90468110 & Depression & --- & Levey et al.\ 2021 \\
\midrule
GCST90468123 & Depression & --- & Levey et al.\ 2021 \\
\midrule
GCST90319327 & Anxiety    & 486,919  & Purves et al.\ 2020 \\
\end{longtable}
}

\subsection*{SI-9.4\quad UniProt Accessions}

{\small
\begin{longtable}{@{}l l l@{}}
\toprule
\textbf{Accession} & \textbf{Gene} & \textbf{Protein} \\
\midrule
\endfirsthead
\toprule
\textbf{Accession} & \textbf{Gene} & \textbf{Protein} \\
\midrule
\endhead
\bottomrule
\endlastfoot
P31645 & \textit{SLC6A4} & Sodium-dependent serotonin transporter \\
\midrule
Q13451 & \textit{FKBP5}  & Peptidyl-prolyl cis-trans isomerase FKBP5 \\
\end{longtable}
}

\subsection*{SI-9.5\quad Reactome Pathway IDs}

{\small
\begin{longtable}{@{}l p{7cm} l@{}}
\toprule
\textbf{Pathway ID} & \textbf{Name} & \textbf{Source Gene} \\
\midrule
\endfirsthead
\toprule
\textbf{Pathway ID} & \textbf{Name} & \textbf{Source Gene} \\
\midrule
\endhead
\bottomrule
\endlastfoot
R-HSA-380615  & Neurotransmitter release cycle --- serotonin & SLC6A4 (P31645) \\
\midrule
R-HSA-442660  & Na$^+$/Cl$^-$ dependent neurotransmitter transporters & SLC6A4 (P31645) \\
\midrule
R-HSA-3371497 & HSP90 chaperone cycle for steroid hormone receptors (SHR) & FKBP5 (Q13451) \\
\end{longtable}
}

\subsection*{SI-9.6\quad KEGG Pathway IDs}

{\small
\begin{longtable}{@{}l l l@{}}
\toprule
\textbf{Pathway ID} & \textbf{Name} & \textbf{Query (Call~\#)} \\
\midrule
\endfirsthead
\toprule
\textbf{Pathway ID} & \textbf{Name} & \textbf{Query (Call~\#)} \\
\midrule
\endhead
\bottomrule
\endlastfoot
map07211 & Serotonin receptor agonists/antagonists & Serotonin (14, 22) \\
\midrule
map04921 & Oxytocin signaling pathway & Oxytocin (18) \\
\midrule
map04927 & Cortisol synthesis and secretion & Cortisol (28) \\
\end{longtable}
}

\subsection*{SI-9.7\quad Human Phenotype Ontology (HPO) Terms}

{\small
\begin{longtable}{@{}l l l@{}}
\toprule
\textbf{HPO ID} & \textbf{Term} & \textbf{Descendants} \\
\midrule
\endfirsthead
\toprule
\textbf{HPO ID} & \textbf{Term} & \textbf{Descendants} \\
\midrule
\endhead
\bottomrule
\endlastfoot
HP:0100851 & Abnormal emotional state & 12 descendants \\
\midrule
HP:0031467 & Dysregulated negative emotional state & 7 descendants \\
\end{longtable}
}

\subsection*{SI-9.8\quad Pharmacology and Pharmacogenomic Identifiers}

\begin{itemize}[nosep]
  \item \textbf{PubChem CID}: 3386 (fluoxetine;
        C$_{17}$H$_{18}$F$_3$NO,
        MW\,=\,309.33\,g/mol)
  \item \textbf{PharmGKB}: PA312 (SLC6A4, chr17q11.2)
  \item \textbf{DGIdb}: 464 interactions across 8 genes
\end{itemize}

\subsection*{SI-9.9\quad WikiPathways}

\begin{itemize}[nosep]
  \item WP5420: ADHD/ASD pathways
  \item WP1455: Serotonin transporter activity
  \item WP727: Monoamine transport
\end{itemize}

\subsection*{SI-9.10\quad Allen Brain Atlas}

\begin{itemize}[nosep]
  \item \textbf{Gene}: Slc6a4 (gene ID 15342)
  \item \textbf{Organisms}: mouse (\textit{Mus musculus}), rat
        (\textit{Rattus norvegicus})
  \item \textbf{Amygdala sub-regions}: 88 anatomical structures
\end{itemize}

\subsection*{SI-9.11\quad Monarch Initiative}

\begin{itemize}[nosep]
  \item HGNC:11050 (SLC6A4)
\end{itemize}

\subsection*{SI-9.12\quad Clinical Trials}

\begin{itemize}[nosep]
  \item \textbf{ClinicalTrials.gov query}: ``emotional dysregulation DBT
        treatment''
  \item \textbf{Total registered trials}: 84
\end{itemize}

\subsection*{SI-9.13\quad Reproducibility Statement}

No private datasets were generated for this review.  All results are
reproducible by re-executing the ToolUniverse tool call sequences documented
in SI-2 using the verbatim JSON arguments provided.

To reproduce the complete analysis:

\begin{enumerate}[nosep]
  \item Install ToolUniverse v1.1.5 and configure MCP stdio mode.
  \item Execute each tool call in SI-2 sequentially, using the verbatim
        JSON arguments provided.
  \item Compare returned results against the ``Key Finding'' column.
  \item Note that API results may vary over time as databases are updated;
        the results reported here reflect the database state at the time
        of query execution (2026).
  \item For STRING enrichment, exact reproduction requires STRING v12 with
        the same background annotation version.  Future STRING versions may
        alter FDR values due to changes in the annotation corpus.
\end{enumerate}
